\chapter{Mở và Khép}
\markboth{Mở và Khép}{Opening and Closing}

\section{Mở lòng để học từ người khác.}
\begin{truyen}{Mở lòng để học từ người khác.}{An open heart learns faster.}
\chuhoa{E}{m từng nghĩ chỉ cần học theo cách của mình là đủ. Nhưng khi chịu mở lòng nghe thêm bạn bè, thầy cô và cả người nhỏ tuổi hơn, em thấy mình hiểu vấn đề nhiều chiều hơn. Mở lòng không làm mất bản sắc; nó cho tri thức thêm cửa đi vào.}
\textit{(The student once thought that learning only in one personal way was enough. But by opening up to classmates, teachers, and even younger voices, understanding became more multi-dimensional. Opening the heart does not erase identity; it creates more doors through which learning can enter.)}
\end{truyen}
\begin{baihoc}
Người mở lòng với điều hay thường trưởng thành nhanh hơn người khép kín trong cái tôi.
\textit{(Those who open themselves to what is good often grow faster than those closed inside ego.)}
\end{baihoc}
\begin{ghinhoanh}
\textbf{Short English to remember:}

Openness invites growth.
\end{ghinhoanh}
\begin{tuhocnhanh}
1. Em có đang khép lại trước góp ý nào có ích không? \textit{(Are you closing yourself off from any useful feedback?)}

2. Em cần mở ra thêm với ai hoặc với điều gì? \textit{(Toward whom or what do you need to open up more?)}
\end{tuhocnhanh}
\ngancach

\section{Khép lại một cánh cửa để dừng bị tổn thương lặp lại.}
\begin{truyen}{Khép lại một cánh cửa để dừng bị tổn thương lặp lại.}{Closing can be an act of self-respect.}
\chuhoa{C}{ô đã cho một mối quan hệ quá nhiều cơ hội, nhưng mỗi lần mở cửa ra là thêm một lần bị xem nhẹ. Đến lúc cô quyết định dừng lại, nhiều người bảo cô lạnh lùng. Thật ra có những cánh cửa nên khép không phải vì hết thương, mà vì còn phải giữ mình.}
\textit{(She gave a relationship too many chances, yet every reopened door led to being undervalued again. When she finally chose to stop, some called her cold. In truth, some doors must be closed not because love is gone, but because self-respect must remain.)}
\end{truyen}
\begin{baihoc}
Mở lòng là đẹp, nhưng không phải mở vô điều kiện cho mọi sự tổn thương quay lại.
\textit{(Openness is beautiful, but it does not mean unconditional access for repeated harm.)}
\end{baihoc}
\begin{ghinhoanh}
\textbf{Short English to remember:}

Some doors close to protect peace.
\end{ghinhoanh}
\begin{tuhocnhanh}
1. Cánh cửa nào em cần khép để giữ bình an? \textit{(Which door do you need to close to protect your peace?)}

2. Em có đang gọi sự tự bảo vệ là lạnh lùng một cách oan uổng không? \textit{(Are you unfairly calling self-protection coldness?)}
\end{tuhocnhanh}
\ngancach

\section{Mở lời xin giúp đỡ thay vì khép trong tự ái.}
\begin{truyen}{Mở lời xin giúp đỡ thay vì khép trong tự ái.}{An open request can break lonely pride.}
\chuhoa{K}{hi bài vở chồng chất, em âm thầm chịu và ngày càng đuối. Chỉ đến khi dám nói với thầy cô và bạn bè rằng mình cần hỗ trợ, em mới thấy mọi thứ sáng ra. Có những người không thiếu người sẵn lòng giúp, chỉ thiếu cánh cửa họ dám mở ra.}
\textit{(When schoolwork piled up, the student endured in silence and became weaker. Only when the courage came to tell teachers and friends that help was needed did the path brighten. Some people do not lack willing helpers; they lack the door they dare to open.)}
\end{truyen}
\begin{baihoc}
Khép trong tự ái làm khó khăn nặng thêm; mở lời đúng lúc làm nó nhẹ đi.
\textit{(Closing yourself within pride makes hardship heavier; asking at the right time makes it lighter.)}
\end{baihoc}
\begin{ghinhoanh}
\textbf{Short English to remember:}

Asking for help is opening, not failing.
\end{ghinhoanh}
\begin{tuhocnhanh}
1. Em có chuyện gì nên nhờ giúp thay vì ôm một mình? \textit{(What should you ask help for instead of carrying it alone?)}

2. Điều gì làm em ngại mở lời? \textit{(What makes you hesitate to ask?)}
\end{tuhocnhanh}
\ngancach

\section{Khép miệng trước chuyện không nên lan truyền.}
\begin{truyen}{Khép miệng trước chuyện không nên lan truyền.}{Closing your mouth can protect another life.}
\chuhoa{M}{ột học sinh biết chuyện riêng của bạn nhưng không mang đi kể như thứ để gây chú ý. Nhờ sự khép miệng ấy, nỗi đau của người kia không biến thành đề tài cho đám đông. Có lúc trưởng thành nằm ở việc biết điều gì nên giữ kín.}
\textit{(A student knew a friend's private struggle but did not spread it as gossip for attention. Because of that closed mouth, the other person's pain did not become public material. Sometimes maturity lies in knowing what should remain protected.)}
\end{truyen}
\begin{baihoc}
Mở lời không phải lúc nào cũng tốt; có điều giữ kín mới là tử tế.
\textit{(Speaking up is not always good; some things are kindest when kept private.)}
\end{baihoc}
\begin{ghinhoanh}
\textbf{Short English to remember:}

Privacy is also a form of care.
\end{ghinhoanh}
\begin{tuhocnhanh}
1. Em có từng kể chuyện người khác chỉ vì thấy thú vị không? \textit{(Have you ever shared someone else's story merely because it seemed interesting?)}

2. Điều gì cần được giữ kín để bảo vệ người liên quan? \textit{(What should remain private in order to protect those involved?)}
\end{tuhocnhanh}
\ngancach

\section{Mở ra với điều mới nhưng biết khép lại với điều độc hại.}
\begin{truyen}{Mở ra với điều mới nhưng biết khép lại với điều độc hại.}{Maturity chooses both welcome and refusal.}
\chuhoa{N}{gười trưởng thành không sống bằng một trạng thái cố định. Họ biết mở ra trước cơ hội học hỏi, tình bạn đẹp, và những góc nhìn lành mạnh; đồng thời cũng biết khép lại với sự xúc phạm, lạm dụng và cám dỗ xấu. Sống khôn là biết cánh cửa nào nên mở, cánh cửa nào nên khóa.}
\textit{(A mature person does not live in one fixed posture. Such a person opens to learning, good friendship, and healthy perspectives, while also knowing how to close against insult, abuse, and destructive temptation. Wise living is knowing which door to open and which to lock.)}
\end{truyen}
\begin{baihoc}
Trưởng thành không phải là mở hết hay khép hết, mà là biết chọn đúng cửa.
\textit{(Maturity is not opening everything or closing everything, but choosing the right door.)}
\end{baihoc}
\begin{ghinhoanh}
\textbf{Short English to remember:}

Wisdom chooses the right door.
\end{ghinhoanh}
\begin{tuhocnhanh}
1. Em đang mở quá mức hay khép quá mức ở mặt nào của đời sống? \textit{(In which area of life are you too open or too closed?)}

2. Một cánh cửa em nên mở và một cánh cửa em nên khép lúc này là gì? \textit{(What is one door you should open and one door you should close right now?)}
\end{tuhocnhanh}
\ngancach
